% GENERATED FILE. Edit canonical lessons, then run npm run generate:books.
% canonical-source-hash: fnv1a64:19e9d45293e16ab7
% canonical-lessons: GE-C143-miete, GE-C143-anzeige, GE-C143-termin, GE-C143-oeffnungszeiten, GE-C143-quittung, GE-R143-first-pass-the-future, GE-R143-second-pass-reading-notices

\chapter{The Rent, an Advert, an Appointment, Opening Hours, a Receipt}
\label{ch:ge-the-rent-an-advert-an-appointment-opening-hours-a-receipt}
\hlchaptermodality{143}

\begin{chapteropening}
\textbf{By the end of this chapter:} \emph{I can read an advert, an appointment, opening hours and a receipt, and act on what they say.}

Complete the last lesson of chapter 143: I can read an advert, an appointment, opening hours and a receipt, and act on what they say.
\end{chapteropening}

\section[die Miete]{die Miete — the rent}

\label{lesson:GE-C143-miete}



\begin{quote}
Before the new one: say the German for next year, then the German for to plan.
\end{quote}

\subsection*{You'll want to know: die Miete}
\textbf{die Miete} — \textquotedblleft{}the rent\textquotedblright{}. The noun behind \emph{mieten}: \textbf{Die Miete ist hoch.} — \textquotedblleft{}The rent is high.\textquotedblright{}

\begin{grammarlens}[title={What you've built}]
One more word for this chapter.
\end{grammarlens}

\subsection*{Your turn}
\begin{itemize}
\raggedright
  \item \emph{Say it:} \emph{die Miete}
  \item \emph{Say it:} \emph{die Miete}, once more
  \item \emph{Say it:} say \emph{die Miete}
  \item \emph{Recall:} say the German for next year, then the German for to plan, then say \emph{die Miete} again
\end{itemize}

\begin{tcolorbox}[breakable,colback=teal!4,colframe=teal!35!black,title={Before you move on}]
What does die Miete mean? (\textquotedblleft{}The rent\textquotedblright{}.) Say it once more.
\end{tcolorbox}
\section[die Anzeige]{die Anzeige — an advert}

\label{lesson:GE-C143-anzeige}



\begin{quote}
Before the new one: say the German for to plan, then the German for the rent.
\end{quote}

\subsection*{You'll want to know: die Anzeige}
\textbf{die Anzeige} — \textquotedblleft{}an advert\textquotedblright{}. Read the advert: \textbf{Zimmer frei. Miete: 400 Euro im Monat.} A room is free, at 400 euros a month.

\begin{grammarlens}[title={What you've built}]
One more word for this chapter.
\end{grammarlens}

\subsection*{Your turn}
\begin{itemize}
\raggedright
  \item \emph{Say it:} \emph{die Anzeige}
  \item \emph{Say it:} \emph{die Anzeige}, once more
  \item \emph{Say it:} say \emph{die Anzeige}
  \item \emph{Recall:} say the German for to plan, then the German for the rent, then say \emph{die Anzeige} again
\end{itemize}

\begin{tcolorbox}[breakable,colback=teal!4,colframe=teal!35!black,title={Before you move on}]
What does die Anzeige mean? (\textquotedblleft{}An advert\textquotedblright{}.) Say it once more.
\end{tcolorbox}
\section[der Termin]{der Termin — an appointment}

\label{lesson:GE-C143-termin}



\begin{quote}
Before the new one: say the German for the rent, then the German for an advert.
\end{quote}

\subsection*{You'll want to know: der Termin}
\textbf{der Termin} — \textquotedblleft{}an appointment\textquotedblright{}. Read the message: \textbf{Ihr Termin: Montag, zehn Uhr.} Your appointment is on Monday at ten: be there before ten.

\begin{grammarlens}[title={What you've built}]
One more word for this chapter.
\end{grammarlens}

\subsection*{Your turn}
\begin{itemize}
\raggedright
  \item \emph{Say it:} \emph{der Termin}
  \item \emph{Say it:} \emph{der Termin}, once more
  \item \emph{Say it:} say \emph{der Termin}
  \item \emph{Recall:} say the German for the rent, then the German for an advert, then say \emph{der Termin} again
\end{itemize}

\begin{tcolorbox}[breakable,colback=teal!4,colframe=teal!35!black,title={Before you move on}]
What does der Termin mean? (\textquotedblleft{}An appointment\textquotedblright{}.) Say it once more.
\end{tcolorbox}
\section[die Öffnungszeiten]{die Öffnungszeiten — opening hours}

\label{lesson:GE-C143-oeffnungszeiten}



\begin{quote}
Before the new one: say the German for an advert, then the German for an appointment.
\end{quote}

\subsection*{You'll want to know: die Öffnungszeiten}
\textbf{die Öffnungszeiten} — \textquotedblleft{}opening hours\textquotedblright{}. Read the door: \textbf{Montag bis Freitag, neun bis achtzehn Uhr.} Open on weekdays from nine to six: not on Saturday.

\begin{grammarlens}[title={What you've built}]
One more word for this chapter.
\end{grammarlens}

\subsection*{Your turn}
\begin{itemize}
\raggedright
  \item \emph{Say it:} \emph{die Öffnungszeiten}
  \item \emph{Say it:} \emph{die Öffnungszeiten}, once more
  \item \emph{Say it:} say \emph{die Öffnungszeiten}
  \item \emph{Recall:} say the German for an advert, then the German for an appointment, then say \emph{die Öffnungszeiten} again
\end{itemize}

\begin{tcolorbox}[breakable,colback=teal!4,colframe=teal!35!black,title={Before you move on}]
What does die Öffnungszeiten mean? (\textquotedblleft{}Opening hours\textquotedblright{}.) Say it once more.
\end{tcolorbox}
\section[die Quittung]{die Quittung — a receipt}

\label{lesson:GE-C143-quittung}



\begin{quote}
Before the new one: say the German for an appointment, then the German for opening hours.
\end{quote}

\subsection*{You'll want to know: die Quittung}
\textbf{die Quittung} — \textquotedblleft{}a receipt\textquotedblright{}. Read the receipt: \textbf{Quittung: Miete, 400 Euro.} The rent is paid: keep it.

\begin{grammarlens}[title={What you've built}]
One more word for this chapter.
\end{grammarlens}

\subsection*{Your turn}
\begin{itemize}
\raggedright
  \item \emph{Say it:} \emph{die Quittung}
  \item \emph{Say it:} \emph{die Quittung}, once more
  \item \emph{Say it:} say \emph{die Quittung}
  \item \emph{Recall:} say the German for an appointment, then the German for opening hours, then say \emph{die Quittung} again
\end{itemize}

\begin{tcolorbox}[breakable,colback=teal!4,colframe=teal!35!black,title={Before you move on}]
What does die Quittung mean? (\textquotedblleft{}A receipt\textquotedblright{}.) Say it once more.
\end{tcolorbox}
\section[(dialogue)]{The future — a first pass}

\label{lesson:GE-R143-first-pass-the-future}



\begin{quote}
Say the words for opening hours and a receipt.
\end{quote}

\subsection*{Your turn}
\begin{itemize}
\raggedright
  \item \emph{Say it:} the future, next week, the day after tomorrow, next year, to plan
\end{itemize}

\begin{tcolorbox}[breakable,colback=teal!4,colframe=teal!35!black,title={Before you move on}]
The future? (\textbf{die Zukunft}.) Next week? (\textbf{nächste Woche}.) The day after tomorrow? (\textbf{übermorgen}.) Next year? (\textbf{nächstes Jahr}.) To plan? (\textbf{vorhaben}.)
\end{tcolorbox}
\section[(dialogue)]{Reading notices — a second pass}

\label{lesson:GE-R143-second-pass-reading-notices}



\begin{quote}
Say the words for the rent and an advert.
\end{quote}

\subsection*{Your turn}
\begin{itemize}
\raggedright
  \item \emph{Say it:} the rent, an advert, an appointment, opening hours, a receipt
\end{itemize}

\begin{tcolorbox}[breakable,colback=teal!4,colframe=teal!35!black,title={Before you move on}]
The rent? (\textbf{die Miete}.) An advert? (\textbf{die Anzeige}.) An appointment? (\textbf{der Termin}.) Opening hours? (\textbf{die Öffnungszeiten}.) A receipt? (\textbf{die Quittung}.)
\end{tcolorbox}
